宽区间和小区间都取301个参数值, 每周期384步, 先积分250周期, 再观察100周期. 宽扫描在$a=.252$首先达到本次人为选择的有限时间指数$.02$阈值, 而$0\le a\le.1$范围的最大估计约$.000489$, 本次观察未找到明确正增长证据. 阈值$.02$、网格间距、初值都是本次实验设置, $.252$不能称系统全局“最早混沌”精确值.

另外用自适应积分舍弃350周期后观察180周期, 复算$a=.2,.252,.05$的指数分别为$.000239,.09536,.000278$. $.2$和$.05$截面接近不变曲线, 在$1,2,4,8,16$次采样移位下均没有小回归误差, 不能将它们误标成有限周期. $.252$有稠密截面和持续正增长的数值证据. 小参数下未观察到混沌, 也不等于证明对全部初值无混沌.

按固定初值扫描未找到可认证的“吸引周期窗口”. 为检查周期结构, 本册另用射击法寻找相同驱动相位下的周期轨道并检验单周期变分矩阵, 结果见下一段. 局部周期岛与同参数其他初值的混沌可以共存, 所以不能把无阻尼模型压缩成一个与初值无关的窗口或第一混沌值.

射击法令相位0时$w(0)=0$, 对$\theta_0$求$w(T_D/2)=0$, 再用驱动的时间反射对称构成整周期轨道. $a=.3$时得到$\theta_0=-.6525804264$, 整周期闭合误差约$4.31\times10^{-11}$. 联立变分方程积分得单周期矩阵行列式$.9999999998$、迹$-1.1680801241$, 两本征值在单位圆上且非实, 是线性椭圆型周期轨道. 另从该点偏移$.002$作180周期采样, 得到围绕周期点的近闭截面曲线. 这给出局部周期轨道及邻近规则运动的数值证据, 与另一初值下的混沌可共存. 椭圆线性分类本身仍不能证明非线性全局稳定. 本题未给出一个对所有初值通用的周期吸引窗口或精确第一混沌参数, 读者复算时应保留此界限.
